\documentclass[10pt,a4paper]{amsart}
\usepackage[margin=19mm]{geometry}
\usepackage{amsmath,amssymb,mathtools}
\usepackage[scr=rsfs,cal=euler]{mathalfa}
\usepackage{xfrac}
\usepackage[unicode,hidelinks,bookmarks=false]{hyperref}
\newcommand{\ie}{i.e.\ }
\newcommand{\cf}{cf.\ }
\newcommand{\resp}{respectively\ }
\newcommand{\Subsection}{\subsection}
\providecommand{\nobreakdash}{\nobreak\hbox{-}}
\setlength{\emergencystretch}{2em}
\multlinegap0pt
\usepackage{bm}
\usepackage{bbm}
\usepackage{dsfont}
\usepackage{xspace}
\usepackage{cancel}
\usepackage{mathtools}
\usepackage{amsthm}
\makeatletter
\@ifundefined{theorem}{\newtheorem{theorem}{Theorem}}{}
\@ifundefined{claim}{\newtheorem{claim}[theorem]{Claim}}{}
\@ifundefined{lemma}{\newtheorem{lemma}[theorem]{Lemma}}{}
\makeatother
\makeatletter
\DeclareMathOperator{\mex}{mex}
\expandafter\ifx\csname proofclaim\endcsname\relax
\newenvironment{proofclaim}{\noindent\emph{Proof of Claim.}}{\qedclaim}
\fi
\newcommand{\xor}{\oplus}
\makeatother
\title[The Normal Play of the Domination Game]{The Normal Play of the Domination Game}
\author{Jo\~ao Marcos Brito, Thiago Marcilon, Nicolas Martins, Rudini Sampaio}
\date{Source: arXiv:2502.13118v4 (CC BY 4.0)}
\begin{document}
\maketitle

\noindent\emph{Source and licence.} The Normal Play of the Domination Game. Version arXiv:2502.13118v4 (CC BY 4.0), distributed under the Creative Commons Attribution 4.0 International licence, as recorded in the arXiv licence metadata for this exact version and shown on the abstract page https://arxiv.org/abs/2502.13118v4. Full rights notice: licenses/arxiv-2502.13118-CC-BY-4.0.txt.\par\smallskip

\noindent\textit{Editorial scope.} Counted unit: the Lemma whose source label is \texttt{lemma1} in the article (source0.tex lines 382--384, source label \texttt{lemma1}), delivered together with the article's own complete original proof of it (source0.tex lines 386--419, one \texttt{proof} environment of 34 lines). No mathematics is written, repaired or re-derived: statement and proof are the article's own text line for line. Required context retained uncounted: claim1 (lines 421--423). Every definition, display and label of those lines is retained unchanged. Omitted from this package and not redistributed: 1--381, 385--385, 420--420, 424--906. Nothing inside a retained excerpt is omitted or altered: every retained body is delivered exactly as the article's lines are written, so no omission receipt of any kind accompanies this package. Citations: the retained excerpts carry no citation command at all, so no bibliography entry of the article is transcribed and no reference is redistributed. Reference adaptation: none was needed and none was made; every reference inside the delivered unit resolves inside it, so no unresolved reference or citation is produced. Printed slips kept literal: no printed word, symbol, hypothesis or number of the counted statement or of its proof was altered. Layout and class: the article's own class is \texttt{article} with packages including amsmath, amssymb, amsthm, xcolor, fullpage, tikz, complexity, enumerate, cite, url, hyperref, pdfpages, subfigure, comment; the wrapper is the lane's pinned \texttt{amsart} editorial wrapper at 10pt on A4, which restates the theorem-like environments the retained text needs and adds the article's own macro definitions that the retained text needs (\texttt{\textbackslash mex}, \texttt{\textbackslash xor}), with extra wrapper packages (bm, bbm, dsfont, xspace, cancel, mathtools). The delivered numbering is the unit's own. Assets: no retained span contains a figure environment, a graphics-inclusion command, a PDF-page inclusion, a table or a TikZ picture; no image file is copied into this package, the package declares no asset dependency and no image file is redistributed with it. Licence: version arXiv:2502.13118v4 of this article is distributed under the Creative Commons Attribution 4.0 International licence, as recorded in the arXiv licence metadata for this exact version and shown on the abstract page https://arxiv.org/abs/2502.13118v4. A full rights notice is delivered at licenses/arxiv-2502.13118-CC-BY-4.0.txt. Characters: every retained excerpt is pure ASCII, so the delivered engine, XeLaTeX with its default Latin Modern text font, reports no missing character.
\begin{claim}\label{claim1}
Let $n$ and $k$ be integers. Then $r(n,k)=(k-3)\% 4\ \xor (n-k)\% 4$ is always different from $n\% 4$, where $\xor$ is the bitwise xor operation. Moreover, $r(n,k)$ is equal to $(n-1)\% 4$ if and only if $n\% 4\in\{1,3\}$ and $k\% 4\in\{0,2\}$.
\end{claim}

\begin{lemma}\label{lemma1}
The nimbers of $P_n'$ and $P_n''$ are equal to $n\% 4$.
\end{lemma}

\begin{proof}
The proof is by induction on $n$.
We first show the four base cases of $P_n''$ with $n=1,2,3,4$.
Since $P_1''$ is a $P_5$ $v_1v_2v_3v_4v_5$ where $v_1$ and $v_5$ were already selected, the three possible moves ($v_2,v_3,v_4$) win immediately (nimber $0$), and then the nimber of $P_1''$ is 1.
Since $P_2''$ is a $P_6$ $v_1\ldots v_6$ where $v_1$ and $v_6$ were already selected, there are four possible moves ($v_2$ to $v_5$), which win immediately (nimber $0$ in $v_3$ or $v_4$) or obtains a $P_1''$ (nimber 1 in $v_2$ or $v_5$), and then the nimber of $P_2''$ is $\mex\{0,1\}=2$.
Since $P_3''$ is a $P_7$ $v_1\ldots v_7$ where $v_1$ and $v_7$ were already selected, there are five possible moves ($v_2$ to $v_6$), which win immediately (nimber $0$ in $v_4$) or obtains a $P_1''$ (nimber 1 in $v_3$ or $v_5$) or obtains a $P_2''$ (nimber 2 in $v_2$ or $v_6$), and then the nimber of $P_3''$ is $\mex\{0,1,2\}=3$.
Finally, since $P_4''$ is a $P_8$ $v_1\ldots v_8$ where $v_1$ and $v_8$ were already selected, there are six possible moves ($v_2$ to $v_7$), which obtains a $P_1''$ (nimber 1 in $v_4$ or $v_5$) or obtains a $P_2''$ (nimber 2 in $v_3$ or $v_6$) or obtains a $P_3''$ (nimber 3 in $v_2$ or $v_7$), and then the nimber of $P_4''$ is $\mex\{1,2,3\}=0$.

Now we show the four base cases of $P_n'$ with $n=1,2,3,4$.
Since $P_1'$ is a $P_3$ $v_1v_2v_3$ where $v_1$ was already selected in the game, the two possible moves ($v_2$ and $v_3$) wins immediately (nimber $0$), and then the nimber of $P_1'$ is 1.
Since $P_2'$ is a $P_4$ $v_1v_2v_3v_4$ where $v_1$ was already selected, there are three possible moves ($v_2,v_3,v_4$), which obtains a $P_1'$ (with nimber 1 in $v_2$) or wins immediately (nimber $0$ in $v_3$ or $v_4$), and then the nimber of $P_2'$ is $\mex\{0,1\}=2$.
Since $P_3'$ is a $P_5$ $v_1v_2v_3v_4v_5$ where $v_1$ was already selected, there are four possible moves ($v_2,v_3,v_4,v_5$), which obtains a $P_2'$ (with nimber 2 in $v_2$) or obtains a $P_1'$ (with nimber 1 in $v_3$) or wins immediately (nimber $0$ in $v_4$) or obtains a $P_1''$ (with nimber 1 in $v_5$), and then the nimber of $P_3'$ is $\mex\{2,1,0,1\}=3$.
Finally, since $P_4'$ is a $P_6$ $v_1v_2v_3v_4v_5v_6$ where $v_1$ was already selected, there are five possible moves ($v_2,v_3,v_4,v_5,v_6$), which obtains a $P_3'$ (with nimber 3 in $v_2$) or obtains a $P_2'$ (with nimber 2 in $v_3$) or obtains a $P_1'$ (with nimber 1 in $v_4$) or obtains a $P_1''$ (with nimber 1 in $v_5$) or obtains a $P_2''$ (with nimber 2 in $v_6$), and then the nimber of $P_4'$ is $\mex\{3,2,1,1,2\}=0$.

So, let $n\geq 5$ and suppose by induction that the lemma holds for every value less than $n$.
First consider the game on $P_n'$ with vertices $v_0,v_1,\ldots,v_{n+1}$, where $v_0$ is the only vertex already selected. Let us enumerate the possible moves of the next player.
If the move is on $v_1$, $v_2$ or $v_3$, it obtains $P_{n-1}'$, $P_{n-2}'$ and $P_{n-3}'$, respectively, with nimbers $(n-1)\% 4$,  $(n-2)\% 4$ and $(n-3)\% 4$ by induction.
If the move is on $v_{n+1}$ or $v_n$, it obtains $P_{n-2}''$ and $P_{n-3}''$, respectively, with nimbers $(n-2)\% 4$ and $(n-3)\% 4$.
So, suppose that the first move is on $v_k$ for $k=4,\ldots,n-1$.
This move obtains two graphs, $P''_{k-3}$ and $P'_{n-k}$, with total nimber $r(n,k)=(k-3)\% 4\ \xor (n-k)\% 4$.
From Claim \ref{claim1} below, $r(n,k)$ is always different from $n\% 4$. Thus, the nimber of $P'_n$ is
\[
\mex\Big\{\ (n-1)\% 4,\ (n-2)\% 4,\ (n-3)\% 4\ \Big\}\ =\ n\% 4.
\]

Now consider the game on $P_n''$ with vertices $v_0,v_1,\ldots,v_{n+3}$, where $v_0$ and $v_{n+3}$ are the only selected vertices. Let us enumerate the possible moves of the next player.
If the move is on $v_1$, $v_2$ or $v_3$, it obtains $P_{n-1}''$, $P_{n-2}''$ and $P_{n-3}''$, respectively, with nimbers $(n-1)\% 4$,  $(n-2)\% 4$ and  $(n-3)\% 4$.
The same if the move is on $v_{n+2}$, $v_{n+1}$ or $v_n$.
So suppose that the move is on $v_k$ for $k=4,\ldots,n-1$.
It obtains two graphs $P''_{k-3}$ and $P''_{n-k}$, with total nimber $r(n,k)=(k-3)\% 4\ \xor (n-k)\% 4$, defined above. From Claim \ref{claim1} below, $r(n,k)$ is always different from $n\% 4$. Thus, the nimber of $P_n''$ is
\[
\mex\Big\{\ (n-1)\% 4,\ (n-2)\% 4,\ (n-3)\% 4\ \Big\}\ =\ n\% 4.
\]
\end{proof}

\end{document}
